\chapter{Methodology}
\label{chapter2}

% ──────────────────────────────────────────────────────────────────────────────────────────────────────────
\section{System Design}

\subsection{Technology Choices}

\paragraph{Programming language.}
Rust was chosen as the implementation language. The primary alternatives considered were C++ and Python. C++
would provide comparable performance but lacks the ownership and borrowing system that prevents common
graphics programming bugs such as dangling pointers into GPU-mapped buffers and use-after-free in callback
closures. Python was ruled out on performance grounds: L-System rewriting and turtle traversal operations are
executed hundreds of thousands of times per scene rebuild, and an interpreted language would make interactive
frame rates unattainable without a native extension layer that would add significant complexity.
Additionally, Rust's type system and trait-based polymorphism allow for a safe and flexible L-System
implementation without sacrificing performance.

\paragraph{Graphics API.}
\texttt{wgpu} was selected as the graphics API. Libraries providing Vulkan and OpenGL bindings such as
\texttt{ash} and \texttt{glow} were considered as alternatives. Whilst targeting Vulkan or OpenGL directly
allows for more control over the rendering process, the higher level abstraction provided by \texttt{wgpu}
reduces the complexity and eliminates platform specific boilerplate. Additionally, as \texttt{wgpu} runs
natively through Vulkan, Metal, DirectX 12 and OpenGL ES~\cite{WgpuPortableGraphics}, using it as the
graphics API increases portability with minimal performance cost.

\paragraph{Windowing and GUI.}
\texttt{winit} provides cross-platform windowing and uses an event-driven model that integrates cleanly with
\texttt{wgpu}'s surface management. For the GUI library \texttt{egui} was chosen over \texttt{imgui} as it is
simpler to integrate with \texttt{wgpu} and its immediate-mode model allows for changes in parameters to be
reflected automatically without the need for explicit event wiring.

\paragraph{Additional dependencies.}
\texttt{glam} provides SIMD-accelerated linear algebra types (\texttt{Vec3}, \texttt{Mat4}) used throughout
the geometry pipeline. \texttt{bytemuck} enables safe zero-copy casting of Rust structs to byte slices for
GPU buffer uploads, with compile-time guarantees of correct memory layout and alignment. The \texttt{rand}
crate provides a seed-controlled \texttt{SmallRng} used for stochastic L-System rules and ecosystem
placement; seeding the generator before each scene build ensures full reproducibility from a single integer
seed. The \texttt{image} crate is used to load the leaf texture atlas from PNG at startup.

% ──────────────────────────────────────────────────────────────────────────────────────────────────────────
\subsection{Architectural Design}
\label{sec:architecture}

The system is structured in five layers as shown in Figure~\ref{fig:arch}. The central design principle is
that each layer communicates through a narrow, well-defined interface, keeping grammar logic, geometry
generation, and rendering fully decoupled from one another.

\begin{figure}[h]
  \centering
  \resizebox{\textwidth}{!}{% High-level system architecture of the Procedural Vegetation Generation system.
% Include in a figure environment with:
%   \resizebox{\textwidth}{!}{% High-level system architecture of the Procedural Vegetation Generation system.
% Include in a figure environment with:
%   \resizebox{\textwidth}{!}{% High-level system architecture of the Procedural Vegetation Generation system.
% Include in a figure environment with:
%   \resizebox{\textwidth}{!}{\input{chapters/figures/architecture}}
\pgfdeclarelayer{background}
\pgfsetlayers{background,main}

\begin{tikzpicture}[
    every node/.style={font=\sffamily\footnotesize},
    % Application / event loop nodes
    app/.style={
      rectangle, rounded corners=3pt,
      draw=blue!70!black, fill=blue!8,
      minimum width=3.2cm, minimum height=0.55cm,
      text centered
    },
    % Settings subsystem nodes
    conf/.style={
      rectangle, rounded corners=3pt,
      draw=green!60!black, fill=green!8,
      minimum width=2.6cm, minimum height=0.55cm,
      text centered
    },
    % World / camera nodes
    world/.style={
      rectangle, rounded corners=3pt,
      draw=orange!80!black, fill=orange!8,
      minimum width=2.6cm, minimum height=0.55cm,
      text centered
    },
    % Scene / controller nodes
    scene/.style={
      rectangle, rounded corners=3pt,
      draw=teal!80!black, fill=teal!8,
      minimum width=3.2cm, minimum height=0.55cm,
      text centered
    },
    % Plant / L-System nodes
    plant/.style={
      rectangle, rounded corners=3pt,
      draw=violet!80!black, fill=violet!8,
      minimum width=2.8cm, minimum height=0.55cm,
      text centered
    },
    % Geometry nodes
    geom/.style={
      rectangle, rounded corners=3pt,
      draw=brown!70!black, fill=brown!8,
      minimum width=2.6cm, minimum height=0.55cm,
      text centered
    },
    % GPU / render nodes
    gpu/.style={
      rectangle, rounded corners=3pt,
      draw=red!70!black, fill=red!8,
      minimum width=3.2cm, minimum height=0.55cm,
      text centered
    },
    % Group box style (parameterised by colour name)
    grp/.style={
      rectangle, rounded corners=6pt,
      draw=#1!60, fill=#1!4,
      inner sep=8pt
    },
    % Primary solid arrow
    arr/.style={-{Stealth[length=5pt]}, thick},
    % Dashed dependency / ownership arrow
    darr/.style={-{Stealth[length=4pt]}, dashed, thin, gray}
  ]

  %% ── Column centres
  % ──────────────────────────────────────────────────────────
  % Left  (Settings):  x =  1.5
  % Centre (App):      x =  6.5
  % Right-mid (Plants):x = 10.5
  % Far-right (World): x = 14.0

  %% ── Application spine
  % ───────────────────────────────────────────────────────
  \node[app] (eventloop) at (0,  0.0) {winit EventLoop};
  \node[app] (app)       at (0, -1.5) {App (ApplicationHandler)};
  \node[app] (state)     at (0, -3.0) {State};

  %% ── Settings subsystem
  % ──────────────────────────────────────────────────────
  \node[conf] (sethead)  at (-5, -3.0) {Settings};
  \node[conf] (setscene) at (-5, -4.0) {SceneSettings};
  \node[conf] (setenv)   at (-5, -5.0) {EnvironmentSettings};
  \node[conf] (setother) at (-5, -6.0) {Camera / Display / LOD};

  \begin{pgfonlayer}{background}
    \node[grp=green, fit=(sethead)(setscene)(setenv)(setother),
    label={[font=\sffamily\tiny,text=green!60!black]above left:settings.rs}]
    (settingsgroup) {};
  \end{pgfonlayer}

  %% ── World subsystem (x=14.0 — clear of the plant group on the
  % right) ────────
  \node[world] (world)  at (5, -3.0) {World};
  \node[world] (camera) at (5, -4.0) {Camera};
  \node[world] (egui)   at (5, -5.0) {EguiRenderer};

  \begin{pgfonlayer}{background}
    \node[grp=orange, fit=(world)(camera)(egui),
    label={[font=\sffamily\tiny,text=orange!70!black]above right:world/}]
    (worldgroup) {};
  \end{pgfonlayer}

  %% ── Renderer
  % ────────────────────────────────────────────────────────────────
  \node[gpu] (renderer) at (0, -4.5) {Renderer};

  %% ── Scene controller
  % ────────────────────────────────────────────────────────
  \node[scene] (scenectrl) at (0, -6.0) {SceneController};

  %% ── Plant / L-System cluster (x=10.5 — clear of world group at x=14.0) ─────
  \node[plant] (ecosystem) at (-4, -8.0) {Ecosystem};
  \node[plant] (plantinst) at (-4, -9.5) {PlantInstance (\texttimes N)};
  \node[plant] (lsystem)   at (-4, -11.0) {LSystem\textless{}S\textgreater{}};
  \node[plant] (rules)     at (-4, -12.5) {Rule\textless{}S\textgreater{}};

  \begin{pgfonlayer}{background}
    \node[grp=violet, fit=(ecosystem)(plantinst)(lsystem)(rules),
    label={[font=\sffamily\tiny,text=violet!70!black]above left:world/plants/}]
    (plantgroup) {};
  \end{pgfonlayer}

  %% ── Turtle / Geometry
  % ───────────────────────────────────────────────────────
  \node[geom] (turtle)  at (4, -8.0) {Turtle};
  \node[geom] (linegeo) at (2, -9.5) {LineGeometry};
  \node[geom] (meshgeo) at (6, -9.5) {MeshGeometry};

  %% ── GPU Buffers + Pipelines
  % ─────────────────────────────────────────────────
  \node[gpu] (gpubuf)       at (4, -11.0) {GPU Buffers (vertex / index)};
  \node[gpu] (pipelined)    at (2, -12.5) {Line Pipeline};
  \node[gpu] (pipelinemesh) at (6, -12.5) {Mesh Pipeline};
  \node[gpu] (display)      at (4, -14.0) {Surface / Display (wgpu)};

  %%─────────────────────────────────────────────────────────────────────────────
  %% ARROWS
  %%─────────────────────────────────────────────────────────────────────────────

  %% Application spine
  \draw[arr, blue!60!black] (eventloop) -- (app)
  node[midway, right, font=\sffamily\tiny] {AppEvent};
  \draw[arr, blue!60!black] (app) -- (state)
  node[midway, right, font=\sffamily\tiny] {AppEvent::Start / input};

  %% State → subsystems (ownership)
  \draw[darr] (state.west) -- +(-0.6,0) |- (sethead.east)
  node[pos=0.3, above, font=\sffamily\tiny, text=gray] {owns};
  \draw[darr] (state.east) -- +(0.6,0) |- (world.west)
  node[pos=0.3, above, font=\sffamily\tiny, text=gray] {owns};

  %% State → Renderer
  \draw[arr, teal!70!black] (state.south) -- (renderer.north)
  node[midway, right, font=\sffamily\tiny] {render()};

  %% Camera → Renderer (view/proj matrix) — stays above the plant group
  \draw[darr] (camera.west) -- +(-1.5,0) |- (renderer.east)
  node[pos=0, above, font=\sffamily\tiny, text=gray] {view\_proj};

  %% Renderer → SceneController
  \draw[arr, teal!70!black] (renderer.south) -- (scenectrl.north)
  node[midway, right, font=\sffamily\tiny] {set\_scene()};

  %% Settings dirty flags → SceneController
  \draw[darr] (setscene.east) -- +(1.5,0) |- (scenectrl.west)
  node[pos=0.3, above, font=\sffamily\tiny, text=gray] {dirty flags};

  %% SceneController → Ecosystem
  \coordinate (sceneunder) at ($(scenectrl.south) + (0, -0.75)$);
  \draw[arr, violet!70!black] (scenectrl.south) -- (sceneunder)
  -| (ecosystem.north);

  %% Ecosystem → PlantInstance → LSystem → Rules
  \draw[arr, violet!70!black] (ecosystem.south) -- (plantinst.north)
  node[midway, right, font=\sffamily\tiny] {per-plant};
  \draw[arr, violet!70!black] (plantinst.south) -- (lsystem.north)
  node[midway, right, font=\sffamily\tiny] {actions()};
  \draw[arr, violet!70!black] (lsystem.south) -- (rules.north)
  node[midway, right, font=\sffamily\tiny] {evolve(n)};

  %% Rules → Turtle: route AROUND the left side to avoid passing through Turtle.
  %% Path: rules.west → waypoint at x=3.0 (clear of all nodes) → turtle.west
  \coordinate (rt) at ($(rules.east) + (2.6, 0)$);
  \draw[arr, brown!70!black] (rules.east)
  -- node[above, font=\sffamily\tiny] {Action[ ]} (rt)
  |- (turtle.west);

  %% Turtle → line/mesh geometry (fork via shared stem)
  \coordinate (turtlestem) at ($(turtle.south) + (0, -0.5)$);
  \draw[brown!70!black, thick] (turtle.south) -- (turtlestem);
  \draw[arr, brown!70!black] (turtlestem) -| (linegeo.north);
  \draw[arr, brown!70!black] (turtlestem) -| (meshgeo.north);

  %% Geometry → GPU buffers
  \coordinate (bufstem) at ($(gpubuf.north) + (0, 0.5)$);
  \draw[arr, red!70!black] (linegeo.south) -- +(0,-0.4) -|
  ($(gpubuf.north) + (-0.6, 0)$);
  \draw[arr, red!70!black] (meshgeo.south) -- +(0,-0.4) -|
  ($(gpubuf.north) + (0.6, 0)$);

  %% GPU buffers → pipelines (fork)
  \coordinate (gpustem) at ($(gpubuf.south) + (0, -0.5)$);
  \draw[red!70!black, thick] (gpubuf.south) -- (gpustem);
  \draw[arr, red!70!black] (gpustem) -| (pipelined.north);
  \draw[arr, red!70!black] (gpustem) -| (pipelinemesh.north);

  %% Pipelines → Display
  \draw[arr, red!70!black] (pipelined.south)    -- +(0,-0.4) -|
  ($(display.north) + (-0.6,0)$);
  \draw[arr, red!70!black] (pipelinemesh.south) -- +(0,-0.4) -|
  ($(display.north) + ( 0.6,0)$);

  %% EguiRenderer → Display: route around the RIGHT edge of the diagram.
  %% Goes east of all nodes, then down, then arrives at display.east.
  \draw[darr] (egui.east)
  -| ($(display.east) + (2,0)$)
  node[pos=0.5, right, font=\sffamily\tiny, text=gray] {UI overlay}
  -| (display.east);

\end{tikzpicture}
}
\pgfdeclarelayer{background}
\pgfsetlayers{background,main}

\begin{tikzpicture}[
    every node/.style={font=\sffamily\footnotesize},
    % Application / event loop nodes
    app/.style={
      rectangle, rounded corners=3pt,
      draw=blue!70!black, fill=blue!8,
      minimum width=3.2cm, minimum height=0.55cm,
      text centered
    },
    % Settings subsystem nodes
    conf/.style={
      rectangle, rounded corners=3pt,
      draw=green!60!black, fill=green!8,
      minimum width=2.6cm, minimum height=0.55cm,
      text centered
    },
    % World / camera nodes
    world/.style={
      rectangle, rounded corners=3pt,
      draw=orange!80!black, fill=orange!8,
      minimum width=2.6cm, minimum height=0.55cm,
      text centered
    },
    % Scene / controller nodes
    scene/.style={
      rectangle, rounded corners=3pt,
      draw=teal!80!black, fill=teal!8,
      minimum width=3.2cm, minimum height=0.55cm,
      text centered
    },
    % Plant / L-System nodes
    plant/.style={
      rectangle, rounded corners=3pt,
      draw=violet!80!black, fill=violet!8,
      minimum width=2.8cm, minimum height=0.55cm,
      text centered
    },
    % Geometry nodes
    geom/.style={
      rectangle, rounded corners=3pt,
      draw=brown!70!black, fill=brown!8,
      minimum width=2.6cm, minimum height=0.55cm,
      text centered
    },
    % GPU / render nodes
    gpu/.style={
      rectangle, rounded corners=3pt,
      draw=red!70!black, fill=red!8,
      minimum width=3.2cm, minimum height=0.55cm,
      text centered
    },
    % Group box style (parameterised by colour name)
    grp/.style={
      rectangle, rounded corners=6pt,
      draw=#1!60, fill=#1!4,
      inner sep=8pt
    },
    % Primary solid arrow
    arr/.style={-{Stealth[length=5pt]}, thick},
    % Dashed dependency / ownership arrow
    darr/.style={-{Stealth[length=4pt]}, dashed, thin, gray}
  ]

  %% ── Column centres
  % ──────────────────────────────────────────────────────────
  % Left  (Settings):  x =  1.5
  % Centre (App):      x =  6.5
  % Right-mid (Plants):x = 10.5
  % Far-right (World): x = 14.0

  %% ── Application spine
  % ───────────────────────────────────────────────────────
  \node[app] (eventloop) at (0,  0.0) {winit EventLoop};
  \node[app] (app)       at (0, -1.5) {App (ApplicationHandler)};
  \node[app] (state)     at (0, -3.0) {State};

  %% ── Settings subsystem
  % ──────────────────────────────────────────────────────
  \node[conf] (sethead)  at (-5, -3.0) {Settings};
  \node[conf] (setscene) at (-5, -4.0) {SceneSettings};
  \node[conf] (setenv)   at (-5, -5.0) {EnvironmentSettings};
  \node[conf] (setother) at (-5, -6.0) {Camera / Display / LOD};

  \begin{pgfonlayer}{background}
    \node[grp=green, fit=(sethead)(setscene)(setenv)(setother),
    label={[font=\sffamily\tiny,text=green!60!black]above left:settings.rs}]
    (settingsgroup) {};
  \end{pgfonlayer}

  %% ── World subsystem (x=14.0 — clear of the plant group on the
  % right) ────────
  \node[world] (world)  at (5, -3.0) {World};
  \node[world] (camera) at (5, -4.0) {Camera};
  \node[world] (egui)   at (5, -5.0) {EguiRenderer};

  \begin{pgfonlayer}{background}
    \node[grp=orange, fit=(world)(camera)(egui),
    label={[font=\sffamily\tiny,text=orange!70!black]above right:world/}]
    (worldgroup) {};
  \end{pgfonlayer}

  %% ── Renderer
  % ────────────────────────────────────────────────────────────────
  \node[gpu] (renderer) at (0, -4.5) {Renderer};

  %% ── Scene controller
  % ────────────────────────────────────────────────────────
  \node[scene] (scenectrl) at (0, -6.0) {SceneController};

  %% ── Plant / L-System cluster (x=10.5 — clear of world group at x=14.0) ─────
  \node[plant] (ecosystem) at (-4, -8.0) {Ecosystem};
  \node[plant] (plantinst) at (-4, -9.5) {PlantInstance (\texttimes N)};
  \node[plant] (lsystem)   at (-4, -11.0) {LSystem\textless{}S\textgreater{}};
  \node[plant] (rules)     at (-4, -12.5) {Rule\textless{}S\textgreater{}};

  \begin{pgfonlayer}{background}
    \node[grp=violet, fit=(ecosystem)(plantinst)(lsystem)(rules),
    label={[font=\sffamily\tiny,text=violet!70!black]above left:world/plants/}]
    (plantgroup) {};
  \end{pgfonlayer}

  %% ── Turtle / Geometry
  % ───────────────────────────────────────────────────────
  \node[geom] (turtle)  at (4, -8.0) {Turtle};
  \node[geom] (linegeo) at (2, -9.5) {LineGeometry};
  \node[geom] (meshgeo) at (6, -9.5) {MeshGeometry};

  %% ── GPU Buffers + Pipelines
  % ─────────────────────────────────────────────────
  \node[gpu] (gpubuf)       at (4, -11.0) {GPU Buffers (vertex / index)};
  \node[gpu] (pipelined)    at (2, -12.5) {Line Pipeline};
  \node[gpu] (pipelinemesh) at (6, -12.5) {Mesh Pipeline};
  \node[gpu] (display)      at (4, -14.0) {Surface / Display (wgpu)};

  %%─────────────────────────────────────────────────────────────────────────────
  %% ARROWS
  %%─────────────────────────────────────────────────────────────────────────────

  %% Application spine
  \draw[arr, blue!60!black] (eventloop) -- (app)
  node[midway, right, font=\sffamily\tiny] {AppEvent};
  \draw[arr, blue!60!black] (app) -- (state)
  node[midway, right, font=\sffamily\tiny] {AppEvent::Start / input};

  %% State → subsystems (ownership)
  \draw[darr] (state.west) -- +(-0.6,0) |- (sethead.east)
  node[pos=0.3, above, font=\sffamily\tiny, text=gray] {owns};
  \draw[darr] (state.east) -- +(0.6,0) |- (world.west)
  node[pos=0.3, above, font=\sffamily\tiny, text=gray] {owns};

  %% State → Renderer
  \draw[arr, teal!70!black] (state.south) -- (renderer.north)
  node[midway, right, font=\sffamily\tiny] {render()};

  %% Camera → Renderer (view/proj matrix) — stays above the plant group
  \draw[darr] (camera.west) -- +(-1.5,0) |- (renderer.east)
  node[pos=0, above, font=\sffamily\tiny, text=gray] {view\_proj};

  %% Renderer → SceneController
  \draw[arr, teal!70!black] (renderer.south) -- (scenectrl.north)
  node[midway, right, font=\sffamily\tiny] {set\_scene()};

  %% Settings dirty flags → SceneController
  \draw[darr] (setscene.east) -- +(1.5,0) |- (scenectrl.west)
  node[pos=0.3, above, font=\sffamily\tiny, text=gray] {dirty flags};

  %% SceneController → Ecosystem
  \coordinate (sceneunder) at ($(scenectrl.south) + (0, -0.75)$);
  \draw[arr, violet!70!black] (scenectrl.south) -- (sceneunder)
  -| (ecosystem.north);

  %% Ecosystem → PlantInstance → LSystem → Rules
  \draw[arr, violet!70!black] (ecosystem.south) -- (plantinst.north)
  node[midway, right, font=\sffamily\tiny] {per-plant};
  \draw[arr, violet!70!black] (plantinst.south) -- (lsystem.north)
  node[midway, right, font=\sffamily\tiny] {actions()};
  \draw[arr, violet!70!black] (lsystem.south) -- (rules.north)
  node[midway, right, font=\sffamily\tiny] {evolve(n)};

  %% Rules → Turtle: route AROUND the left side to avoid passing through Turtle.
  %% Path: rules.west → waypoint at x=3.0 (clear of all nodes) → turtle.west
  \coordinate (rt) at ($(rules.east) + (2.6, 0)$);
  \draw[arr, brown!70!black] (rules.east)
  -- node[above, font=\sffamily\tiny] {Action[ ]} (rt)
  |- (turtle.west);

  %% Turtle → line/mesh geometry (fork via shared stem)
  \coordinate (turtlestem) at ($(turtle.south) + (0, -0.5)$);
  \draw[brown!70!black, thick] (turtle.south) -- (turtlestem);
  \draw[arr, brown!70!black] (turtlestem) -| (linegeo.north);
  \draw[arr, brown!70!black] (turtlestem) -| (meshgeo.north);

  %% Geometry → GPU buffers
  \coordinate (bufstem) at ($(gpubuf.north) + (0, 0.5)$);
  \draw[arr, red!70!black] (linegeo.south) -- +(0,-0.4) -|
  ($(gpubuf.north) + (-0.6, 0)$);
  \draw[arr, red!70!black] (meshgeo.south) -- +(0,-0.4) -|
  ($(gpubuf.north) + (0.6, 0)$);

  %% GPU buffers → pipelines (fork)
  \coordinate (gpustem) at ($(gpubuf.south) + (0, -0.5)$);
  \draw[red!70!black, thick] (gpubuf.south) -- (gpustem);
  \draw[arr, red!70!black] (gpustem) -| (pipelined.north);
  \draw[arr, red!70!black] (gpustem) -| (pipelinemesh.north);

  %% Pipelines → Display
  \draw[arr, red!70!black] (pipelined.south)    -- +(0,-0.4) -|
  ($(display.north) + (-0.6,0)$);
  \draw[arr, red!70!black] (pipelinemesh.south) -- +(0,-0.4) -|
  ($(display.north) + ( 0.6,0)$);

  %% EguiRenderer → Display: route around the RIGHT edge of the diagram.
  %% Goes east of all nodes, then down, then arrives at display.east.
  \draw[darr] (egui.east)
  -| ($(display.east) + (2,0)$)
  node[pos=0.5, right, font=\sffamily\tiny, text=gray] {UI overlay}
  -| (display.east);

\end{tikzpicture}
}
\pgfdeclarelayer{background}
\pgfsetlayers{background,main}

\begin{tikzpicture}[
    every node/.style={font=\sffamily\footnotesize},
    % Application / event loop nodes
    app/.style={
      rectangle, rounded corners=3pt,
      draw=blue!70!black, fill=blue!8,
      minimum width=3.2cm, minimum height=0.55cm,
      text centered
    },
    % Settings subsystem nodes
    conf/.style={
      rectangle, rounded corners=3pt,
      draw=green!60!black, fill=green!8,
      minimum width=2.6cm, minimum height=0.55cm,
      text centered
    },
    % World / camera nodes
    world/.style={
      rectangle, rounded corners=3pt,
      draw=orange!80!black, fill=orange!8,
      minimum width=2.6cm, minimum height=0.55cm,
      text centered
    },
    % Scene / controller nodes
    scene/.style={
      rectangle, rounded corners=3pt,
      draw=teal!80!black, fill=teal!8,
      minimum width=3.2cm, minimum height=0.55cm,
      text centered
    },
    % Plant / L-System nodes
    plant/.style={
      rectangle, rounded corners=3pt,
      draw=violet!80!black, fill=violet!8,
      minimum width=2.8cm, minimum height=0.55cm,
      text centered
    },
    % Geometry nodes
    geom/.style={
      rectangle, rounded corners=3pt,
      draw=brown!70!black, fill=brown!8,
      minimum width=2.6cm, minimum height=0.55cm,
      text centered
    },
    % GPU / render nodes
    gpu/.style={
      rectangle, rounded corners=3pt,
      draw=red!70!black, fill=red!8,
      minimum width=3.2cm, minimum height=0.55cm,
      text centered
    },
    % Group box style (parameterised by colour name)
    grp/.style={
      rectangle, rounded corners=6pt,
      draw=#1!60, fill=#1!4,
      inner sep=8pt
    },
    % Primary solid arrow
    arr/.style={-{Stealth[length=5pt]}, thick},
    % Dashed dependency / ownership arrow
    darr/.style={-{Stealth[length=4pt]}, dashed, thin, gray}
  ]

  %% ── Column centres
  % ──────────────────────────────────────────────────────────
  % Left  (Settings):  x =  1.5
  % Centre (App):      x =  6.5
  % Right-mid (Plants):x = 10.5
  % Far-right (World): x = 14.0

  %% ── Application spine
  % ───────────────────────────────────────────────────────
  \node[app] (eventloop) at (0,  0.0) {winit EventLoop};
  \node[app] (app)       at (0, -1.5) {App (ApplicationHandler)};
  \node[app] (state)     at (0, -3.0) {State};

  %% ── Settings subsystem
  % ──────────────────────────────────────────────────────
  \node[conf] (sethead)  at (-5, -3.0) {Settings};
  \node[conf] (setscene) at (-5, -4.0) {SceneSettings};
  \node[conf] (setenv)   at (-5, -5.0) {EnvironmentSettings};
  \node[conf] (setother) at (-5, -6.0) {Camera / Display / LOD};

  \begin{pgfonlayer}{background}
    \node[grp=green, fit=(sethead)(setscene)(setenv)(setother),
    label={[font=\sffamily\tiny,text=green!60!black]above left:settings.rs}]
    (settingsgroup) {};
  \end{pgfonlayer}

  %% ── World subsystem (x=14.0 — clear of the plant group on the
  % right) ────────
  \node[world] (world)  at (5, -3.0) {World};
  \node[world] (camera) at (5, -4.0) {Camera};
  \node[world] (egui)   at (5, -5.0) {EguiRenderer};

  \begin{pgfonlayer}{background}
    \node[grp=orange, fit=(world)(camera)(egui),
    label={[font=\sffamily\tiny,text=orange!70!black]above right:world/}]
    (worldgroup) {};
  \end{pgfonlayer}

  %% ── Renderer
  % ────────────────────────────────────────────────────────────────
  \node[gpu] (renderer) at (0, -4.5) {Renderer};

  %% ── Scene controller
  % ────────────────────────────────────────────────────────
  \node[scene] (scenectrl) at (0, -6.0) {SceneController};

  %% ── Plant / L-System cluster (x=10.5 — clear of world group at x=14.0) ─────
  \node[plant] (ecosystem) at (-4, -8.0) {Ecosystem};
  \node[plant] (plantinst) at (-4, -9.5) {PlantInstance (\texttimes N)};
  \node[plant] (lsystem)   at (-4, -11.0) {LSystem\textless{}S\textgreater{}};
  \node[plant] (rules)     at (-4, -12.5) {Rule\textless{}S\textgreater{}};

  \begin{pgfonlayer}{background}
    \node[grp=violet, fit=(ecosystem)(plantinst)(lsystem)(rules),
    label={[font=\sffamily\tiny,text=violet!70!black]above left:world/plants/}]
    (plantgroup) {};
  \end{pgfonlayer}

  %% ── Turtle / Geometry
  % ───────────────────────────────────────────────────────
  \node[geom] (turtle)  at (4, -8.0) {Turtle};
  \node[geom] (linegeo) at (2, -9.5) {LineGeometry};
  \node[geom] (meshgeo) at (6, -9.5) {MeshGeometry};

  %% ── GPU Buffers + Pipelines
  % ─────────────────────────────────────────────────
  \node[gpu] (gpubuf)       at (4, -11.0) {GPU Buffers (vertex / index)};
  \node[gpu] (pipelined)    at (2, -12.5) {Line Pipeline};
  \node[gpu] (pipelinemesh) at (6, -12.5) {Mesh Pipeline};
  \node[gpu] (display)      at (4, -14.0) {Surface / Display (wgpu)};

  %%─────────────────────────────────────────────────────────────────────────────
  %% ARROWS
  %%─────────────────────────────────────────────────────────────────────────────

  %% Application spine
  \draw[arr, blue!60!black] (eventloop) -- (app)
  node[midway, right, font=\sffamily\tiny] {AppEvent};
  \draw[arr, blue!60!black] (app) -- (state)
  node[midway, right, font=\sffamily\tiny] {AppEvent::Start / input};

  %% State → subsystems (ownership)
  \draw[darr] (state.west) -- +(-0.6,0) |- (sethead.east)
  node[pos=0.3, above, font=\sffamily\tiny, text=gray] {owns};
  \draw[darr] (state.east) -- +(0.6,0) |- (world.west)
  node[pos=0.3, above, font=\sffamily\tiny, text=gray] {owns};

  %% State → Renderer
  \draw[arr, teal!70!black] (state.south) -- (renderer.north)
  node[midway, right, font=\sffamily\tiny] {render()};

  %% Camera → Renderer (view/proj matrix) — stays above the plant group
  \draw[darr] (camera.west) -- +(-1.5,0) |- (renderer.east)
  node[pos=0, above, font=\sffamily\tiny, text=gray] {view\_proj};

  %% Renderer → SceneController
  \draw[arr, teal!70!black] (renderer.south) -- (scenectrl.north)
  node[midway, right, font=\sffamily\tiny] {set\_scene()};

  %% Settings dirty flags → SceneController
  \draw[darr] (setscene.east) -- +(1.5,0) |- (scenectrl.west)
  node[pos=0.3, above, font=\sffamily\tiny, text=gray] {dirty flags};

  %% SceneController → Ecosystem
  \coordinate (sceneunder) at ($(scenectrl.south) + (0, -0.75)$);
  \draw[arr, violet!70!black] (scenectrl.south) -- (sceneunder)
  -| (ecosystem.north);

  %% Ecosystem → PlantInstance → LSystem → Rules
  \draw[arr, violet!70!black] (ecosystem.south) -- (plantinst.north)
  node[midway, right, font=\sffamily\tiny] {per-plant};
  \draw[arr, violet!70!black] (plantinst.south) -- (lsystem.north)
  node[midway, right, font=\sffamily\tiny] {actions()};
  \draw[arr, violet!70!black] (lsystem.south) -- (rules.north)
  node[midway, right, font=\sffamily\tiny] {evolve(n)};

  %% Rules → Turtle: route AROUND the left side to avoid passing through Turtle.
  %% Path: rules.west → waypoint at x=3.0 (clear of all nodes) → turtle.west
  \coordinate (rt) at ($(rules.east) + (2.6, 0)$);
  \draw[arr, brown!70!black] (rules.east)
  -- node[above, font=\sffamily\tiny] {Action[ ]} (rt)
  |- (turtle.west);

  %% Turtle → line/mesh geometry (fork via shared stem)
  \coordinate (turtlestem) at ($(turtle.south) + (0, -0.5)$);
  \draw[brown!70!black, thick] (turtle.south) -- (turtlestem);
  \draw[arr, brown!70!black] (turtlestem) -| (linegeo.north);
  \draw[arr, brown!70!black] (turtlestem) -| (meshgeo.north);

  %% Geometry → GPU buffers
  \coordinate (bufstem) at ($(gpubuf.north) + (0, 0.5)$);
  \draw[arr, red!70!black] (linegeo.south) -- +(0,-0.4) -|
  ($(gpubuf.north) + (-0.6, 0)$);
  \draw[arr, red!70!black] (meshgeo.south) -- +(0,-0.4) -|
  ($(gpubuf.north) + (0.6, 0)$);

  %% GPU buffers → pipelines (fork)
  \coordinate (gpustem) at ($(gpubuf.south) + (0, -0.5)$);
  \draw[red!70!black, thick] (gpubuf.south) -- (gpustem);
  \draw[arr, red!70!black] (gpustem) -| (pipelined.north);
  \draw[arr, red!70!black] (gpustem) -| (pipelinemesh.north);

  %% Pipelines → Display
  \draw[arr, red!70!black] (pipelined.south)    -- +(0,-0.4) -|
  ($(display.north) + (-0.6,0)$);
  \draw[arr, red!70!black] (pipelinemesh.south) -- +(0,-0.4) -|
  ($(display.north) + ( 0.6,0)$);

  %% EguiRenderer → Display: route around the RIGHT edge of the diagram.
  %% Goes east of all nodes, then down, then arrives at display.east.
  \draw[darr] (egui.east)
  -| ($(display.east) + (2,0)$)
  node[pos=0.5, right, font=\sffamily\tiny, text=gray] {UI overlay}
  -| (display.east);

\end{tikzpicture}
}
  \caption{System architecture. Solid arrows denote call/data flow; dashed arrows denote ownership or
    configuration dependencies. The \texttt{SceneController} is the interface between the rendering layer
  and the plant/L-System subsystem.}
  \label{fig:arch}
\end{figure}

\paragraph{Application layer.}
The \texttt{State} struct owns the \texttt{EguiRenderer}, \texttt{World}, and \texttt{Settings} and drives
the \texttt{winit} event loop via the \texttt{ApplicationHandler} trait. Keyboard shortcuts for scene cycling
and parameter stepping are handled here, keeping input logic out of the rendering path. Settings are
organised into sub-structs (\texttt{SceneSettings}, \texttt{EnvironmentSettings}, \texttt{CameraSettings},
\texttt{DisplaySettings}, \texttt{LodSettings}, and \texttt{CullSettings}) each carrying a generation counter
incremented by \texttt{mark\_dirty()}; the scene controller rebuilds only when the observed generation
differs from its last-seen value.

\paragraph{World and scene control.}
\texttt{World} owns the camera, projection, and \texttt{SceneController}. The \texttt{SceneController} calls
\texttt{actions()} on each plant, passes the returned action sequence through a \texttt{Turtle} instance, and
uploads the resulting geometry to pre-allocated GPU buffers. LOD tier is computed per-plant from camera
distance at build time; tier changes also mark the scene dirty, triggering a rebuild at the next update.
Ecosystem regeneration uses the current seed from \texttt{EnvironmentSettings}, so every rebuild from the
same seed produces the same spatial layout.

\paragraph{Plant layer.}
Each species is a struct implementing the \texttt{Plant} trait. The L-System string is cached and only
re-evolved when the iteration count changes or the season crosses a threshold, so grammar computation does
not occur every frame. Seasonal behaviour, dormancy, shedding probability, leaf colouration, and flowering
phenology, is encoded directly within the parametric production closures, which capture season-derived values
at generate time.

\paragraph{Turtle and geometry.}
The \texttt{Turtle} struct maintains its position, an orthonormal frame (heading, normal, binormal) and a
matrix stack for push/pop. It produces two geometry streams: \texttt{LineGeometry} for the debug skeleton and
\texttt{MeshGeometry} for cylinder branches and leaf quads. Tropism, gravitropism, wind, voxel-grid space
pruning, and tapered branch radii are applied per-segment during traversal.

\paragraph{Renderer.}
Two \texttt{wgpu} render pipelines are used: a line pipeline for the debug branch skeleton and a mesh
pipeline for cylinder branches and leaf quads. The mesh pipeline accepts position, colour, normal, UV, and
tangent vertex buffers and a 96-byte uniform block (view-projection matrix, light direction, ambient factor,
and season colour tint). The season tint is written to the uniform buffer every frame, so seasonal colour
changes require no geometry rebuild.

% ──────────────────────────────────────────────────────────────────────────────────────────────────────────
\section{Development}

\subsection{Development Approach}

An incremental, minimum viable product (MVP) strategy was used, motivated by the risk that later components
(ecosystem placement, full seasonal L-System effects) might prove infeasible within the project timescale.
Development proceeded in priority order, ensuring a visually functioning system existed at every stage. This
approach mirrors the bottom-up ordering used by Deussen et al.~\cite{deussenRealisticModelingRendering1998},
where a working single-plant renderer was established before any ecosystem-level features were attempted.

\medskip
All code is version-controlled with Git on a single main branch with frequent commits. Additionally, a nix
flake was used to manage the projects' development environment and provide declarative, reproducible builds
across machines~\cite{dolstraPurelyFunctionalSoftware2006, nixManual}.

\subsection{Development Sprints}

Development was organised into four sprints, each targeting a subset of the project deliverables from
Chapter~\ref{chapter1}, with a defined goal and deliverabilities assessed at the end of the sprint.

\subsubsection{Sprint~1 — L-System engine and 2D rendering}
\textit{Goal:} A working L-System interpreter capable of producing correct string derivations, interpreted by
a turtle into a 2D line drawing rendered via wgpu.

\medskip
\textit{Deliverables:}
\begin{itemize}
  \item The \texttt{LSystem} struct, \texttt{Symbol} trait, and \texttt{Rule} enum supporting \texttt{Normal}
    and \texttt{Stochastic}, \texttt{ContextSensitive} and \texttt{Parametric} variants (Deliverable~2)
  \item A 2D turtle producing \texttt{LineGeometry}
  \item A \texttt{wgpu} pipeline capable of rendering \texttt{LineGeometry}
\end{itemize}

\subsubsection{Sprint~2 — 3D geometry and plant grammars}
\textit{Goal:} A 3D turtle producing cylinder branches and leaf quads, with at least three distinct plant
grammars demonstrating different grammar classes.

\medskip
\textit{Deliverables:}
\begin{itemize}
  \item The 3D turtle with push/pop state, branch, and leaf quad generation producing \texttt{MeshGeometry}
  \item A \texttt{wgpu} pipeline capable of rendering \texttt{MeshGeometry} (Deliverable~1)
  \item The \texttt{Plant} trait and \texttt{SceneController}
  \item Three distinct plant grammars (Deliverable~3)
\end{itemize}

\subsubsection{Sprint~3 — Environmental response and GUI}
\textit{Goal:} Phototropism, gravitropism, wind, seasonal changes and space constraints with a GUI to manage
these parameters.

\medskip
\textit{Deliverables:}
\begin{itemize}
  \item Implementation of phototropism, gravitropism and wind response
  \item Implementation of Seasonal dormancy, colouration, shedding and flowering phenology
  \item A GUI from controlling environmental parameters (Deliverable~4)
\end{itemize}

\subsubsection{Sprint~4 — Ecosystem placement and LOD}
\textit{Goal:} A kernel-based ecosystem placement algorithm with self-thinning and succession as well as a
distance-based LOD system.

\medskip
\textit{Deliverables:}
\begin{itemize}
  \item Inhibitory, promotional, mixed, and neutral placement kernels with inverse-transform sampling
  \item Implementation of self-thinning and succession
  \item A five-tier LOD system with adaptive geometry budget pressure
  \item View-frustum culling using the Gribb--Hartmann method~\cite{gribbFastExtractionViewing}
\end{itemize}
